\documentclass[11pt,a4paper]{article}
			\usepackage[french]{babel}
					
				\usepackage{pifont}  
				\usepackage[utf8x]{inputenc}
				\usepackage[T1]{fontenc} 
				\usepackage{lmodern}			
				\usepackage{fancyhdr}
				\usepackage{textcomp}
				\usepackage{makeidx}
				\usepackage{tabularx}
				\usepackage{multicol}
				\usepackage{multirow}
				\usepackage{longtable}
				\usepackage{color}
				\usepackage{soul}
				\usepackage{boxedminipage}
				\usepackage{shadow}
				\usepackage{framed}			
				\usepackage{array}
				\usepackage{url}
				\usepackage{ragged2e}
				\usepackage{fancybox}
				\newcommand{\cadretitre}[2]{
				  \vspace*{0.8\baselineskip}
				  \begin{center}%
				  \boxput*(0,1){%
					%\colorbox{white}{\Large\textbf{\ #1\ }}%
				  }%
				  {%
					\setlength{\fboxsep}{10pt}%
				    \Ovalbox{\begin{minipage}{.8\linewidth}\begin{center}\Large\sffamily{#2}\end{center}\end{minipage}}}%
				  \end{center}
				  \vspace*{2\baselineskip}
				  }
			
			\makeatletter
			\def\@seccntformat#1{\protect\makebox[0pt][r]{\csname the#1\endcsname\quad}}
			\makeatother

				% Permet d'afficher qqchose à une positin absolue
				\usepackage[absolute]{textpos}
				\setlength{\TPHorizModule}{1cm}
				\setlength{\TPVertModule}{\TPHorizModule}
	
				\usepackage[titles]{tocloft}
				\setlength{\cftbeforesecskip}{0.5ex}
				\setlength{\cftbeforesubsecskip}{0.2ex}
				\addto\captionsfrench{\renewcommand\contentsname{}}
				
				\usepackage[font=scriptsize]{caption}
				
				\usepackage{listings}
\lstdefinestyle{lstverb}
  {
    basicstyle=\footnotesize,
    frameround=tttt, frame=trbl, framerule=0pt, rulecolor=\color{gray},
    lineskip=-1pt,   % pour rapprocher les lignes
    flexiblecolumns, escapechar=\\,
    tabsize=4, extendedchars=true
  }
\lstnewenvironment{Java}[1][]{\lstset{style=lstverb,language=java,#1}}{}
				\ifx\pdfoutput\undefined
					\usepackage{graphicx}
				\else
					\usepackage[pdftex]{graphicx}
				\fi
				\usepackage[a4paper, hyperfigures=true, colorlinks, linkcolor=black, citecolor=blue,urlcolor=blue, pagebackref=true, bookmarks=true, bookmarksopen=true,bookmarksnumbered=true,
                pdfauthor={}, pdftitle={TD Variables structurées}, pdfkeywords={TD Variables structur\'ees, },pdfpagemode=UseOutlines,pdfpagetransition=Dissolve,nesting=true,
				backref, pdffitwindow=true, bookmarksnumbered=true]{hyperref}
				\usepackage{supertabular}
				\usepackage[table]{xcolor}
				\usepackage{url}
				\usepackage{caption} 
				\setlength{\parskip}{1.3ex plus 0.2ex minus 0.2ex}
				\setlength{\parindent}{0pt}
				
				\makeatletter
				\def\url@leostyle{ \@ifundefined{selectfont}{\def\UrlFont{\sf}}{\def\UrlFont{\footnotesize\ttfamily}}}
				\makeatother
				\urlstyle{leo}
				
				\definecolor{examplecolor}{rgb}{0.156,0.333,0.443}
				\definecolor{definitioncolor}{rgb}{0.709,0.784,0.454}
				\definecolor{exercisecolor}{rgb}{0.49,0.639,0}
				\definecolor{hintcolor}{rgb}{0.941,0.674,0.196}
				\definecolor{tableHeadercolor}{rgb}{0.709,0.784,0.454}
				\definecolor{tablerowAltcolor}{rgb}{.866,.905,.737}
				\definecolor{tablerowAlt2color}{rgb}{.968,.976,.933}
				\definecolor{verylightgray}{rgb}{0.98,0.98,0.98}
				
				\newenvironment{fshaded}{
				\def\FrameCommand{\fcolorbox{framecolor}{shadecolor}}
				\MakeFramed {\FrameRestore}}
				{\endMakeFramed}
				
				\newenvironment{fexample}[1][]{\definecolor{shadecolor}{rgb}{.913,.913,.913}
				\definecolor{framecolor}{rgb}{.156,.333,.443}
				\begin{fshaded}}{\end{fshaded}} 
				
				\newenvironment{fdefinition}{\definecolor{shadecolor}{rgb}{.913,.913,.913}
				\definecolor{framecolor}{rgb}{.709,.784,.454}
				\begin{fshaded}}{\end{fshaded}}
				
				\newenvironment{fexercise}{\definecolor{shadecolor}{rgb}{.913,.913,.913}
				\definecolor{framecolor}{rgb}{.49,.639,0}
				\begin{fshaded}}{\end{fshaded}}
				
				\newenvironment{fhint}{\definecolor{shadecolor}{rgb}{.913,.913,.913}
				\definecolor{framecolor}{rgb}{.941,.674,.196}
				\begin{fshaded}}{\end{fshaded}}	
				
				\newcommand{\PreserveBackslash}[1]{
				\let\temp=\\#1\let\\=\temp
				}
				\let\PBS=\PreserveBackslash
				\newcolumntype{A}{>{\PBS\raggedright\small\hspace{0pt}}X}
				\newcolumntype{L}[1]{>{\PBS\raggedright\small\hspace{0pt}}p{#1}}
				\newcolumntype{R}[1]{>{\PBS\raggedleft\small\hspace{0pt}}p{#1}}
				\newcolumntype{C}[1]{>{\PBS\centering\small\hspace{0pt}}p{#1}}
				
				\makeindex
				
				\title{TD Variables structur\'ees}	
			\date{}
			\author{\scriptsize{}}
			\definecolor{light-gray}{gray}{0.8}
			\renewcommand{\headrulewidth}{0pt}
			\fancyhead[L]{
				\footnotesize\textsc{Haute \'Ecole de Bruxelles}\\
	    			\footnotesize\textsc{\'Ecole Sup\'erieure d'Informatique}
			}
			\fancyhead[R]{
				\footnotesize{Bachelor en Informatique}\\
				\footnotesize{Laboratoires Java} - 
			\footnotesize{1\`ere ann\'ee}}
				\fancyfoot[L]{ }
				\fancyfoot[C]{}
				\fancyfoot[R]{\scriptsize{\textcolor{gray}{version 2014-2015 (\today)}}}
				\pagestyle{plain}
				\reversemarginpar
				\usepackage{rotating}						
				\begin{document}
					\begin{textblock}{9}(2,3.2)
						\includegraphics[width=2cm]{../../../_templates/java/icons/logo-esi}
					\end{textblock}
				
				
				
				
				%\maketitle
				\cadretitre{TD1}{TD Variables structur\'ees}
				\thispagestyle{fancy}
        \marginpar{\begin{sideways}
            \begin{minipage}[t]{1cm}
            \begin{tiny}
            \includegraphics[width=1\linewidth,height=1\textheight,keepaspectratio=true]{../../../_templates/java/icons/cc-gris.jpg}
			\end{tiny}
			\end{minipage}
            \begin{minipage}[b]{19cm}
            \begin{tiny}
            \textcolor{gray}{Distribué sous licence Creative Commons Paternité - Partage à l'Identique 2.0 Belgique 
            (\texttt{http://creativecommons.org/licenses/by-sa/2.0/be/})
			\vspace{-1em}
			\\Les autorisations au-delà du champ de cette licence peuvent être obtenues à 
			\texttt{http://www.heb.be/esi}
			- \texttt{mcodutti@heb.be}
			}\end{tiny}
			\end{minipage}
        \end{sideways}}
            \begin{abstract}
			Ce TD a pour but de rappeler les variables structur\'ees.
		
            \par
        \end{abstract}
				\vspace{-2em}\tableofcontents
				\pagestyle{plain}
            \clearpage
            \fancyhead[L,C,R]{}
            \fancyfoot[L,C]{}
            \fancyfoot[R]{ \scriptsize{\textcolor{gray}{
				TDStruct - page \thepage}}}
				\thispagestyle{fancy}
				\pagestyle{fancy}
	   
            \section{Algorithmes s\'equentiels}
        Dans les chapitres pr\'ec\'edents, nous avons utilis\'e des variables de types dits \guillemotleft  simples \guillemotright 
        (entier, r\'eel, bool\'een, caract\`ere, chaine) ne pouvant contenir qu'une seule valeur \`a la fois.
        Cependant, certains types d'information consistent en un regroupement de plusieurs donn\'ees
        \'el\'ementaires. Quelques exemples :
        
					\begin{itemize}
				
			\item Une date est compos\'ee de trois \'el\'ements (le jour, le mois, l'ann\'ee). Le mois peut
            s'exprimer par un entier (15/10/2014) ou par une chaine (15 octobre 2014).
          
			\item 
            Un moment de la journ\'ee est un triple d'entiers (heures, minutes, secondes).
          
			\item 
            La localisation d'un point dans un plan n\'ecessite la connaissance de deux coordonn\'ees
            cart\'esiennes (l'abscisse x et l'ordonn\'ee y) ou polaires (le rayon et l'angle).
          
			\item 
            Une adresse est compos\'ee de plusieurs donn\'ees : un nom de rue, un num\'ero de maison,
            parfois un num\'ero de boite postale, un code postal, le nom de la localit\'e, un nom ou
            code de pays pour une adresse \`a l'\'etranger. . .
          
					\end{itemize}
				
            \par
        
        Pour stocker et manipuler de tels ensembles de donn\'ees, nous utiliserons des 
        \textbf{types structur\'es} ou \textbf{structures}. 
        Une structure est donc un \textbf{ensemble fini d'\'el\'ements pouvant \^etre de types distincts}. 
        Chaque \'el\'ement de cet ensemble, appel\'e \textbf{champ} de la structure, poss\`ede un nom unique.
      
            \par
        
        Notez qu'un champ d'une structure peut lui-m\^eme \^etre une structure. Par exemple, une
        carte d'identit\'e inclut parmi ses informations une date de naissance, l'adresse de son
        propri\'etaire (cach\'ee dans la puce \'electronique !). . .
      
            \par
        \subsection{D\'efinition d'une structure}
		    La d\'efinition d'un type structur\'e adoptera le mod\`ele suivant :
      
            \par
        \begin{verbatim}
structure NomDeLaStructure
    nomChamp1 : type1
    nomChamp2 : type2
    ...
    nomChampN : typeN
fin structure
      \end{verbatim}\verb@nomChamp1@, . . ., \verb@nomChampN@ 
        sont les noms des diff\'erents champs de la structure, et\verb@ type1@,
        . . ., \verb@typeN@ les types de ces champs. 
        Ces types sont soit les types \guillemotleft  simples \guillemotright  \'etudi\'es pr\'ec\'edemment 
        (entier, r\'eel, bool\'een, caract\`ere, chaine) soit d'autres types structur\'es dont la
        structure aura \'et\'e pr\'ealablement d\'efinie.
      
            \par
        
        Pour exemple, nous d\'efinissons ci-dessous trois types structur\'es que nous utiliserons souvent par la suite :
      
            \par
        \begin{verbatim}
structure Date
    jour : entier
    mois : entier
    année : entier
fin structure
      \end{verbatim}\begin{verbatim}
structure Moment
    heure : entier
    minute : entier
    seconde : entier
fin structure
      \end{verbatim}\begin{verbatim}
structure Point
    x : réel
    y : réel
fin structure
      \end{verbatim}\subsection{D\'eclaration d'une variable de type structur\'e}
		    Une fois un type structur\'e d\'efini, la d\'eclaration de variables de ce type est identique \`a celle des variables simples. Par exemple :
      
            \par
        \begin{verbatim}
anniversaire, jourJ : Date
départ, arrivée, unMoment : Moment
a, b, centreGravité : Point
      \end{verbatim}\subsection{Utilisation des variables de type structur\'e}
		    En g\'en\'eral, pour manipuler une variable structur\'ee ou en modifier le contenu, il faut agir
        au niveau de ses champs en utilisant les op\'erations permises selon leur type. Pour acc\'eder \`a
        l'un des champs d'une variable structur\'ee, il faut mentionner le nom de ce champ ainsi que
        celui de la variable dont il fait partie. Nous utiliserons la notation \guillemotleft  point\'ee \guillemotright  :
      
            \par
        \begin{verbatim}
nomVariable.nomChamp
      \end{verbatim}
		    Exemple d'instructions utilisant les variables d\'eclar\'ees au paragraphe pr\'ec\'edent :
      
            \par
        \begin{verbatim}
anniversaire.jour ← 15
anniversaire.mois ← 10
anniversaire.année ← 2014
arrivée.heure ← départ.heure + 2
centreGravité.x ← (a.x + b.x) / 2
      \end{verbatim}
		    On peut cependant, dans certains cas, utiliser une variable structur\'ee de mani\`ere globale
        (c'est-\`a-dire d'une fa\c con qui agit simultan\'ement sur chacun de ses champs). Le cas le plus
        courant est l'affectation interne entre deux variables structur\'ees de m\^eme type, par exemple :
      
            \par
        \begin{verbatim}
arrivée ← départ
      \end{verbatim}
		    qui r\'esume les trois instructions suivantes :
      
            \par
        \begin{verbatim}
arrivée.heure ← départ.heure
arrivée.minute ← départ.minute
arrivée.seconde ← départ.seconde
      \end{verbatim}
		    Une variable structur\'ee peut aussi \^etre le param\`etre d'un module, et un module peut \'egalement 
		    renvoyer une \guillemotleft  valeur \guillemotright  de type structur\'e. Par exemple, l'ent\^ete d'un module renvoyant
		    le nombre de secondes \'ecoul\'ees entre une heure de d\'epart et d'arriv\'ee serait :
      
            \par
        \begin{verbatim}
module nbSecondesEcoulées( départ↓, arrivée↓ : Moment) → entier
      \end{verbatim}
		    On pourra aussi lire ou afficher une variable structur\'ee :
      
            \par
        \begin{verbatim}
lire unMoment
afficher unMoment
      \end{verbatim}
		    Par contre, il n'est pas autoris\'e d'utiliser les op\'erateurs de comparaison (<, >) pour comparer 
		    des variables structur\'ees (m\^eme de m\^eme type), car une relation d'ordre n'accompagne
        pas toujours les structures utilis\'ees. En effet, s'il est naturel de vouloir comparer des dates
        ou des moments, comment d\'efinir une relation d'ordre avec les points du plan ou avec des
        cartes d'identit\'es ?
      
            \par
        
        Si le besoin de comparer des variables structur\'ees se fait sentir, il faudra dans ce cas \'ecrire
        des modules de comparaison adapt\'es aux structures utilis\'ees.
      
            \par
        
        Par facilit\'e d'\'ecriture, on peut assigner tous les champs en une fois via des \guillemotleft  {} \guillemotright . Exemple :
      
            \par
        \begin{verbatim}
anniversaire ← {1, 9, 1989}
      \end{verbatim}\subsection{Exemple d'algorithme}
		    Le module ci-dessous re\c coit en param\`etre deux dates ; la valeur renvoy\'ee est -1 si la premi\`ere
        date est ant\'erieure \`a la deuxi\`eme, 0 si les deux dates sont identiques ou 1 si la premi\`ere date
        est post\'erieure \`a la deuxi\`eme.
      
            \par
        \begin{verbatim}
// Reçoit 2 dates en paramètres et retourne la valeur
// –1 si la première date est antérieure à la deuxième
// 0 si les deux dates sont identiques
// 1 si la première date est postérieure à la deuxième.
module comparerDates(date1↓, date2↓ : Date) → entier
    résultat : entier
    résultat ← -1 // valeur choisie par défaut
    si date1.année ≥ date2.année alors
      si date1.année > date2.année alors
        résultat ← 1
      sinon // les années sont identiques
        si date1.mois ≥ date2.mois alors
          si date1.mois > date2.mois alors
            résultat ← 1
          sinon // les années et les mois sont identiques
            si date1.jour ≥ date2.jour alors
              si date1.jour > date2.jour alors
                résultat ← 1
              sinon // les années, les mois et les jours sont identiques
                résultat ← 0
              fin si
            fin si
          fin si
        fin si
      fin si
    fin si
    retourner résultat
fin module
      \end{verbatim}\section{Exercices}
				Maintenant, mettons tout \c ca en pratique.
      
            \par
        \subsection{\`A vous de jouer...}
        \'Ecrivez un algorithme qui 
          
					\begin{enumerate}
				
			\item re\c coit en param\`etre un moment d'une journ\'ee et qui retourne le nombre de secondes \'ecoul\'ees depuis minuit jusqu'\`a ce moment.
			\item 
              re\c coit en param\`etre un nombre de secondes \'ecoul\'ees dans une journ\'ee 
              \`a partir de minuit et qui retourne le moment correspondant de la journ\'ee.
			\item re\c coit en param\`etres deux moments d'une journ\'ee et qui retourne le
              nombre de secondes s\'eparant ces deux moments.
			\item recevant les coordonn\'ees de deux points distincts du plan et qui retourne
              les coordonn\'ees du point situ\'e au milieu des deux.
			\item recevant les coordonn\'ees de deux points distincts du plan et qui retourne
              la distance entre ces deux points.
					\end{enumerate}
				
            \par
        
			
		\subparagraph{Cercle} 
		
					\textcolor{white}{.} \par
				
          D\'efinir un type Cercle pouvant d\'ecrire de fa\c con commode un cercle quelconque dans un
          espace \`a deux dimensions. \'Ecrire ensuite
          
					\begin{itemize}
				
			\item un module calculant la surface du cercle re\c cu en param\`etre ;
			\item un module recevant 2 points en param\`etre et retournant le cercle dont le diam\`etre est le segment reliant ces 2 points ;
			\item un module qui d\'etermine si un point donn\'e est dans un cercle ;
			\item un module qui indique si 2 cercles ont une intersection.
					\end{itemize}
				
            \par
        
				\end{document}
			